% Katalog G: Formelsatz und Zielketten (Palette Slot 7).
\ZSFNewColumn
\StartChapter[kat:formeln]{G · Formeln}

\begin{runintext}
  Die Formel-API ist immer aktiv; die erweiterten Operatoren stehen in
  \ZSFref{kat:module}. Jeder Eintrag zeigt die Makros in ihrer gesetzten
  Form, nicht ihren Aufruf.
\end{runintext}

\SubsectionBar[kat:formel-basis]{Basis-Makros}
\ZSFindex{Zahlmenge}
\ZSFindex{Begrenzer}

\begin{formulabox}[G-01 · Zahlmengen]
  \[
    x\in\R,\quad z\in\C,\quad n\in\N,\quad k\in\Z,\quad q\in\Q
  \]
\end{formulabox}
\Zweck{R, C, N, Z, Q statt rohem mathbb.}

\begin{formulabox}[G-02 · Operatoren und Begrenzer]
  \[
    \dd x,\quad \abs{x},\quad \norm{\vect v},\quad \sgn(x)
  \]
  \ZSFsep
  \[
    \grad f,\quad \divg \vect F,\quad \rot \vect F
  \]
\end{formulabox}
\Zweck{Aufrechtes Differential, gepaarte Begrenzer, Vektorsatz und die drei
Feldoperatoren.}

\begin{formulabox}[G-03 · Summe und Limes]
  \[
    \ZSFsumAuto{i=1}{n} i,\quad \ZSFlimAuto{x\to0}\frac{\sin x}{x}
  \]
\end{formulabox}
\Zweck{Der Grenzen-Modus (side, stack, auto) wird zentral umgeschaltet, nicht
pro Formel.}

\SubsectionBar[kat:formel-anno]{Anmerkung und Benennung}
\ZSFindex{formulanote}
\ZSFindex{ZSFbraceunder}

\begin{formulabox}[G-04 · Klammern]
  \[
    \ZSFbraceunder{x^2+2x+1}{Ausgangsform}
    = \ZSFbraceover{(x+1)^2}{faktorisiert}
  \]
\end{formulabox}
\Zweck{Der Name steht am Teil statt daneben — für einen zusammenhängenden
Term der Normalfall.}

\begin{formulabox}[G-05 · Anmerkungen]
  \ZSFformulaline{a^2+b^2=c^2}{G-06 · ZSFformulaline}
  \ZSFsep
  \[
    E = mc^2
  \]
  \formulanote{G-07 · formulanote als eigene Zeile darunter.}
\end{formulabox}
\Zweck{Drei Rollen, ein Stil: rechtsbündig auf der Zeile, als eigene Zeile
darunter, oder als ganzer Satz über textNachBox.}

\SubsectionBar[kat:formel-farbe]{Terme verfolgen}
\ZSFindex{ZSFmhlA}

\begin{formulabox}[G-08 · ZSFmhlA bis D]
  \[
    \ZSFmhlA{a}x^2 + \ZSFmhlB{b}x + \ZSFmhlD{c} = 0
  \]
  \ZSFsep
  \[
    \ZSFmhlC{x_{1,2}} = \frac{-\ZSFmhlB{b} \pm
      \sqrt{\ZSFmhlB{b}^2 - 4\,\ZSFmhlA{a}\,\ZSFmhlD{c}}}{2\,\ZSFmhlA{a}}
  \]
\end{formulabox}
\Zweck{Positional: A, B und D sind die Stränge, C ist immer das Ziel. Gilt
innerhalb einer Herleitung — das Gegenstück sind die Grössenfarben (E-16).}

\SubsectionBar[kat:formel-goal]{Zielketten}
\ZSFindex{goalbox}
\ZSFindex{GoalCondition}

\begin{goalbox}[G-09 · Primitive]
  \GoalCondition{Bedingung}
  \GoalStep{\text{Zwischenschritt}}
  \GoalStep{\text{beliebig oft wiederholbar}}
  \GoalTarget{\text{Ziel}}
\end{goalbox}
\Zweck{Frei gebaute Kette: eine Bedingung, beliebig viele Schritte, ein Ziel.}

\begin{goalbox}[G-10 · ZSFDerivationCase]
  \ZSFDerivationCase*{Sternform, einzeilig}{\sqrt{x^2}}{x}
  \ZSFsep
  \ZSFDerivationCase{Grundform}{\sqrt{x^2}}{-x}
\end{goalbox}
\Zweck{Der fertige Dreisatz Bedingung, Ausdruck, Resultat für die
Fallunterscheidung.}
